\scanpage{103}
\hrulefill

Soit $X/K$, et soit $T$ une partie de
\inserted{\tikz[baseline=0pt,x=1em,y=1ex]{
  \node[inner sep=0pt] at (0,1.8) {$X=$};
  \draw[line width=.4pt] (-.65,1.0) .. controls (-.3,.45) and (.45,.55) .. (.45,-.1) -- (.62,-1.8);
}}$X_{\mathbb C}(\mathbb C)$, loc.\textsuperscript{t} fermée 1-connexe. Soit $X'$ \deleted{une variété algébrique}\inserted{$-X$ un} $K$-morphisme, qui sur \deleted{la partie de} un voisinage de $T$ soit fini étale. Alors $X'\vert T=X'_T$ se décompose en un nb fini de copies de $T$ -- il y a une relation d'équivalence dans $X'_T$, qu'on appelle la $T$-adjacence de points de $X'$ au dessus de $T$. On se donne \deleted{$I\subset X'_T$} $x',y'\in X'_T$
% The following small baseline trace is distinct from the overwork on the inserted T.
\unsure{\tikz[baseline=0pt,x=1ex,y=1ex]{
  \draw[line width=.45pt] (.6,.8) -- (.32,.58)
    .. controls (-.05,.55) and (-.18,.2) .. (.05,.08)
    .. controls (.38,-.02) and (.4,.3) .. (.28,.44)
    .. controls (.35,-.1) and (.15,-.5) .. (.05,-.7);
}}
\inserted{$T$ adjacents;}
\margin{\unsure{en supp.\textsuperscript{t}|en supposant} $x',y'\in X'_{\overline K}$}
si $u\in\Gamma$, on ne peut affirmer que \deleted{\unsure{ils le soient}}\inserted{$u(x'),u(y')$} soient encore $T$ adjacents.
Mais s'ils le sont, on peut regarder $l^T_{y',x'}:x'\to y'$ \inserted{``$T$-horizontal''} dans \deleted{$\unsure{\in}\,\unsure{\widehat{\Pi}_1(X'_{\overline K},x',y')}$} $\operatorname{Hom}_{\widehat{\Pi}_1(X'_{\overline K})}(x',y')$ et son transformé, \deleted{\unsure{le} \illegible{1} \unsure{de} $T$ qui passe par $x',y'$} \deleted{et son transformé} et le \uline{comparer} \deleted{$\unsure{u(l_{y',x'})}$}. : $l^T_{u(y'),u(x')}$ (qui est défini par hyp. d'horizontalité)
% The period belongs physically just after the cancelled lower formula;
% the separate colon is higher, after comparer.
\[
(53)\qquad \deleted{\unsure{u}}{}^u l^T_{y',x'}=g^T_{x,y}(u)\,l^T_{u(y'),u(x')}
\]
où
\[
g^{T\deleted{\unsure{u}}}_{x,y}(u)\in\widehat{\pi}_1\bigl(X',
\underset{\substack{\text{strate horizontale}\\\text{de }X'_T\text{ qui contient}\\u(x'),u(y')}}{
  \tikz[baseline=(term.base)]{
    \node[inner sep=0pt] (term) {$T(u x',u y')$};
    \draw[line width=.4pt]
      ([yshift=-1pt]term.south west) -- ([yshift=-4pt]term.south west)
      -- ([yshift=-4pt]term.south east) -- ([yshift=-1pt]term.south east);
  }
}\bigr)
\]
Ceci posé, on aura
\[
(54)\qquad \widehat f\underbrace{\bigl(g^T_{y',x'}(u)\bigr)}_{\substack{\rotatebox{-90}{$\in$}\\\widehat{\pi}_1(X,T)}}=g^T_{y,x}(u)
\]
\mspage{102}
